\chapter{Capstone 学期收尾与下一届 Warm Start}

\section{案例导读：学期收尾要为下一届保留起跑线}

Capstone 的结束不是简单打包提交。真正好的收尾会把本届学生的项目经验、失败样本、修订建议、优秀示例和维护提醒整理出来，让下一届不是从空白开始。

本章把 shutdown 和 warm start 放在一起，是为了提醒课程团队：结束和开始其实是同一个闭环的两端。只有把本学期的证据整理好，下一轮课程才能更快进入有效学习。

\section{案例目标}

第 55 章把临时替班、应急接管和 minimal run package 整理成 substitute handoff workflow。本章继续处理一轮课程真正结束时的另一面：\textbf{怎样在学期结束后完成 shutdown checklist，同时给下一届学生和下一轮课程团队留下一个可直接启动的 warm-start package。}

本章的目标有四点：

\begin{enumerate}
    \item 理解 end-of-term shutdown 不是“学期结束了就归档”，而是把当前轮次安全关机并给下一轮留下可接续的种子；
    \item 学会检查 shutdown check、archive snapshot、next-cohort seed、owner 和 release pressure；
    \item 学会区分“可以收尾”“补 shutdown 检查”“先补归档快照”“准备下一届种子包”和“先指定收尾 owner”五类出口；
    \item 学会把一轮课程的结束变成下一轮课程的更好起点。
\end{enumerate}

\section{学期结束不是简单打包}

课程运行到最后，团队通常最累，也最容易把“终于结束了”误当成“已经处理完了”。但真正影响下一轮课程质量的，往往不是本学期最后一次展示，而是收尾阶段有没有把 submission、feedback、reproducibility note、TA digest 和 course retrospective 处理清楚。

如果 shutdown 做得太匆忙，下一轮团队会拿到一堆文件，但不知道哪些是最终版、哪些只是临时导出、哪些值得继续用、哪些应该淘汰。更糟的是，下一届学生可能面对一份表面完整、实际上没有 warm-start 入口的课程包。

所以，本章的核心立场是：\textbf{课程收尾的目标不是结束，而是为下一轮留下可用、可懂、可接续的起始状态。}

\section{一个极小的 shutdown-warmstart 数据集}

配套 notebook 使用 \nolinkurl{capstone_shutdown_warmstart_demo.csv}。这是一组完全本地、教学用途的\textbf{学期收尾与下一届预热卡片}。每一行代表一个学期末需要收尾、归档或留种的模块，包含：

\begin{itemize}
    \item \texttt{shutdown\_id}；
    \item \texttt{shutdown\_area}；
    \item \texttt{carry\_forward\_scope}；
    \item \texttt{shutdown\_check\_status}；
    \item \texttt{archive\_snapshot\_status}；
    \item \texttt{next\_cohort\_seed\_status}；
    \item \texttt{owner\_status}；
    \item \texttt{release\_pressure\_score}；
    \item \texttt{reference\_route}。
\end{itemize}

这里的 route 分成五类：

\begin{itemize}
    \item \texttt{shutdown\_ready}：可以进入学期收尾包；
    \item \texttt{complete\_shutdown\_check}：还有关机前检查没做完；
    \item \texttt{archive\_before\_shutdown}：归档快照、版本或证据还不完整；
    \item \texttt{prepare\_next\_cohort\_seed}：还没有把内容整理成下一届可接手的起始包；
    \item \texttt{assign\_shutdown\_owner}：没有明确的收尾负责人。
\end{itemize}

\section{先看一个危险 baseline：只按收尾压力排序}

学期末很容易按 release pressure 排序：越赶、越怕忘掉的事情越先做。

\begin{examplebox}[只按 release pressure 选择收尾任务]
\begin{lstlisting}[style=pythonstyle]
top4 = sorted(
    rows,
    key=lambda row: row["release_pressure_score"],
    reverse=True,
)[:4]
\end{lstlisting}
\end{examplebox}

这个 baseline 的问题在于，压力高不一定代表该任务已经 ready。它可能还缺 shutdown check、archive snapshot，或者根本没有整理成下一届可以直接接手的种子包。

在本章教学数据上，只按收尾压力取前四项，真正可收尾的精度只有 0.25，未就绪项混入率是 0.75。表~\ref{tab:ch56_pressure_vs_shutdown} 展示了结构化 shutdown workflow 如何先看 owner，再看 shutdown check，再看 archive snapshot，最后检查 next-cohort seed：

\begin{table}[htbp]
\centering
\small
\setlength{\tabcolsep}{4pt}
\begin{tabular}{@{}p{0.28\textwidth}>{\centering\arraybackslash}p{0.18\textwidth}>{\centering\arraybackslash}p{0.18\textwidth}p{0.28\textwidth}@{}}
\toprule
方法 & 可收尾精度 & 未就绪混入率 & 教学解读 \\
\midrule
只按收尾压力取前四项 & 0.25 & 0.75 & 紧迫感不能替代完整检查、归档和 warm-start 准备 \\
结构化 shutdown workflow & 1.00 & 0.00 & 先定人、再关机、再留档、最后留种 \\
\bottomrule
\end{tabular}
\caption{本章 notebook 中两个最小学期收尾流程的对照。课程收尾不是压力清单，而是过渡状态表。}
\label{tab:ch56_pressure_vs_shutdown}
\end{table}

\section{一个可执行的 shutdown workflow}

本章 notebook 的 routing 逻辑可以写成：

\begin{examplebox}[一个最小学期收尾 routing 逻辑]
\begin{lstlisting}[style=pythonstyle]
if owner_status != "assigned":
    assign_shutdown_owner
elif shutdown_check_status != "complete":
    complete_shutdown_check
elif archive_snapshot_status != "complete":
    archive_before_shutdown
elif next_cohort_seed_status != "seeded":
    prepare_next_cohort_seed
else:
    shutdown_ready
\end{lstlisting}
\end{examplebox}

这个顺序体现了一个过渡原则：\textbf{负责人先于收尾，关机检查先于归档，归档先于 warm start。}

\section{学期收尾包应该包含什么}

一份可执行的 end-of-term shutdown package 至少应该包含：

\begin{enumerate}
    \item 最终版 submission 与评分边界说明；
    \item 学生与助教反馈摘要；
    \item 可复现的课程快照：notebook、数据、关键路径和版本；
    \item course retrospective：哪些地方最耗时、最易错、最值得保留；
    \item 下一届 warm-start 种子：starter repo、calendar seed、最小 handout 和 kickoff note；
    \item owner 与最后检查日期：谁关机、谁归档、谁给下一轮开机。
\end{enumerate}

它不需要保存一切，但必须留下“下一轮从哪里开始”这个最重要的问题的答案。

\section{配套 notebook 做了什么}

本章 notebook 采用的是一个 teaching-first 的最小设计：

\begin{itemize}
    \item 先读取 shutdown-warmstart table，并统计 carry-forward scope、shutdown check、archive、seed 和 reference route；
    \item 用 release pressure 做一个 naive shutdown baseline；
    \item 再实现一个带 owner、shutdown check、archive snapshot 和 next-cohort seed 的 shutdown workflow；
    \item 输出 shutdown-ready、complete-check、archive-first、prepare-seed 和 assign-owner 五个队列；
    \item 为每个非 ready 模块生成收尾前 checklist；
    \item 最后生成 shutdown package outline 和 warm-start assistant prompt。
\end{itemize}

\section{AI 助手如何帮助你}

在这个案例里，AI 助手最适合帮助你：

\begin{itemize}
    \item 把一学期的杂乱材料压缩成 shutdown package；
    \item 从反馈和课程记录里提炼出最值得给下一届保留的种子；
    \item 标出哪些内容适合 carry forward、哪些应该淘汰；
    \item 为下一届 kickoff 生成一页 warm-start briefing；
    \item 起草 course retrospective 和 next-cohort starter note。
\end{itemize}

但 AI 助手不能替团队决定哪些学生材料可以长期保留，也不能替课程负责人完成最终归档判断。它能让收尾更清楚，却不能替你结束这一轮课程。

\section{练习}

\begin{exercisebox}[基础练习]
把一个 \texttt{shutdown\_ready} 模块的 \texttt{next\_cohort\_seed\_status} 改成 \texttt{missing}。重新运行 notebook，观察它为什么会离开 ready 队列。
\end{exercisebox}

\begin{exercisebox}[进阶练习]
新增一个 release pressure 很高、但 \texttt{archive\_snapshot\_status = missing} 的模块。预测它会落到哪个 route，并说明为什么“着急结束”不等于“可以收尾”。
\end{exercisebox}

\begin{exercisebox}[开放练习]
为你自己的课程项目写一页 end-of-term shutdown brief。至少包含：最终版交付、feedback digest、reproducibility snapshot、course retrospective、next-cohort seed 和 owner。
\end{exercisebox}

\section{本章小结}

一轮课程的结束不应该只是归档，而应该是下一轮课程的更好起点。本章把 shutdown check、archive snapshot、next-cohort seed、owner 和 release pressure 放进同一张 shutdown card。

当课程团队能把模块分成可以收尾、补 shutdown 检查、先补归档快照、准备下一届种子包和先指定收尾 owner，Capstone 就从一学期的工作堆积变成可以跨届延续的教学系统。
